\chapter{2000年京皖卷（春文）}

\section{单选题}

\begin{problem}
复数 \(z_1=3+\i\)，\(z_2=1-\i\)，则 \(z=z_1z_2\) 在复平面内的对应点位于（\(\quad\)）

\choices
    {第一象限}
    {第二象限}
    {第三象限}
    {第四象限}
\end{problem}

\begin{answer}
D
\end{answer}

\begin{solution}
由复数乘法得 \(z=(3+\i)(1-\i)=3-3\i+\i-\i^2=4-2\i\). 因为 \(z\) 的实部为正、虚部为负，所以它在复平面的第四象限，故选 D.
\end{solution}

\begin{problem}
设全集 \(I=\{a,b,c,d,e\}\)，集合 \(M=\{a,c,d\}\)，\(N=\{b,d,e\}\)，那么 \(\overline{M}\cap\overline{N}\) 是（\(\quad\)）

\choices
    {\(\varnothing\)}
    {\(\{d\}\)}
    {\(\{a,c\}\)}
    {\(\{b,e\}\)}
\end{problem}

\begin{answer}
A
\end{answer}

\begin{solution}
全集为 \(I\)，所以 \(\overline{M}=\{b,e\}\)，\(\overline{N}=\{a,c\}\). 两者没有公共元素，即 \(\overline{M}\cap\overline{N}=\varnothing\)，故选 A.
\end{solution}

\begin{problem}
已知双曲线 \(\frac{x^2}{b^2} - \frac{y^2}{a^2} = 1\) 的两条渐近线互相垂直，那么该双曲线的离心率是（\(\quad\)）

\choices
    {\(2\)}
    {\(\sqrt{3}\)}
    {\(\sqrt{2}\)}
    {\(\frac{3}{2}\)}
\end{problem}

\begin{answer}
C
\end{answer}

\begin{solution}
双曲线的两条渐近线为 \(y=\pm\frac{a}{b}x\)，其斜率之积为 \(-\frac{a^2}{b^2}\). 两渐近线垂直，故 \(-\frac{a^2}{b^2}=-1\)，从而 \(a=b\). 该双曲线的焦半距满足 \(c^2=a^2+b^2=2b^2\)，因此离心率 \(e=\frac{c}{b}=\sqrt{2}\)，故选 C.
\end{solution}

\begin{problem}
下列方程关于 \(x = y\) 对称的是（\(\quad\)）

\choices
    {\(x^{2} - x + y^{2} = 1\)}
    {\(x^{2}y + xy^{2} = 1\)}
    {\(x - y = 1\)}
    {\(x^{2} - y^{2} = 1\)}
\end{problem}

\begin{answer}
B
\end{answer}

\begin{solution}
关于直线 \(x=y\) 对称的充要条件是交换 \(x\)，\(y\) 后方程保持不变. 选项 A 交换后为 \(y^2-y+x^2=1\)，与原方程不同；选项 B 交换后仍为 \(y^2x+yx^2=1\)，即 \(x^2y+xy^2=1\)，方程不变；选项 C 交换后为 \(y-x=1\)，即 \(x-y=-1\)，与原方程不同；选项 D 交换后为 \(y^2-x^2=1\)，即 \(x^2-y^2=-1\)，与原方程不同. 因此选 B.
\end{solution}

\begin{problem}
一个圆锥的底面直径和高都同一个球的直径相等，那么圆锥与球的体积之比是（\(\quad\)）

\choices
    {\(1:3\)}
    {\(2:3\)}
    {\(1:2\)}
    {\(2:9\)}
\end{problem}

\begin{answer}
C
\end{answer}

\begin{solution}
设球的半径为 \(r\)，则球的直径为 \(2r\). 圆锥的底面半径为 \(r\)，高为 \(2r\)，所以 \(V_{\text{锥}}=\frac13\pi r^2\cdot2r=\frac23\pi r^3\)，而 \(V_{\text{球}}=\frac43\pi r^3\). 因此 \(V_{\text{锥}}:V_{\text{球}}=1:2\)，故选 C.
\end{solution}

\begin{problem}
直线 \((\sqrt{3} - \sqrt{2})x + y = 3\) 和直线 \(x + (\sqrt{2} - \sqrt{3})y = 2\) 的位置关系是（\(\quad\)）

\choices
    {相交但不垂直}
    {垂直}
    {平行}
    {重合}
\end{problem}

\begin{answer}
B
\end{answer}

\begin{solution}
第一条直线的斜率为 \(m_1=\sqrt{2}-\sqrt{3}\)，第二条直线的斜率为 \(m_2=-\frac{1}{\sqrt{2}-\sqrt{3}}=\sqrt{2}+\sqrt{3}\). 于是 \(m_1m_2=(\sqrt{2}-\sqrt{3})(\sqrt{2}+\sqrt{3})=-1\)，两直线垂直，故选 B.
\end{solution}

\begin{problem}
函数 \(y=\lg|x|\)（\(\quad\)）

\choices
    {是偶函数，在区间 \((-\infty,0)\) 上单调递增}
    {是偶函数，在区间 \((-\infty,0)\) 上单调递减}
    {是奇函数，在区间 \((0,+\infty)\) 上单调递增}
    {是奇函数，在区间 \((0,+\infty)\) 上单调递减}
\end{problem}

\begin{answer}
B
\end{answer}

\begin{solution}
函数的定义域为 \(x\neq0\)，且 \(f(-x)=\lg|-x|=\lg|x|=f(x)\)，所以它是偶函数. 在 \(x<0\) 时，\(f'(x)=\frac{1}{x\ln10}<0\)，故函数在 \((-\infty,0)\) 上单调递减，选 B.
\end{solution}

\begin{problem}
从单词“equation”选取 \(5\) 个不同的字母排成一排，含有“qu”（其中“qu”相连且顺序不变）的不同排列共有（\(\quad\)）

\choices
    {\(120\) 个}
    {\(480\) 个}
    {\(720\) 个}
    {\(840\) 个}
\end{problem}

\begin{answer}
B
\end{answer}

\begin{solution}
“equation”中的 \(8\) 个字母各不相同. 除 \(q\)，\(u\) 外，从其余 \(6\) 个字母中选 \(3\) 个，有 \(\mathrm C_6^3\) 种；把相连且顺序固定的“qu”看作一个整体，与另外 \(3\) 个字母共 \(4\) 个对象排列，有 \(4!\) 种. 因此总数为 \(\mathrm C_6^3\cdot4!=20\cdot24=480\)，故选 B.
\end{solution}

\begin{problem}
椭圆短轴长是 \(2\)，长轴长是短轴的 \(2\) 倍，则椭圆中心到其准线距离是（\(\quad\)）

\choices
    {\(\frac{8}{5}\sqrt{5}\)}
    {\(\frac{4}{5}\sqrt{5}\)}
    {\(\frac{8}{3}\sqrt{3}\)}
    {\(\frac{4}{3}\sqrt{3}\)}
\end{problem}

\begin{answer}
D
\end{answer}

\begin{solution}
椭圆短轴长为 \(2\)，故短半轴 \(b=1\)；长轴长为短轴的 \(2\) 倍，故长半轴 \(a=2\). 焦半距 \(c=\sqrt{a^2-b^2}=\sqrt{3}\)，离心率 \(e=\frac ca\). 椭圆中心到准线的距离为 \(\frac ae=\frac{a^2}{c}=\frac4{\sqrt3}=\frac43\sqrt3\)，故选 D.
\end{solution}

\begin{problem}
函数 \(y=\sin x+\cos x+2\) 的最小值是（\(\quad\)）

\choices
    {\(2 - \sqrt{2}\)}
    {\(2 + \sqrt{2}\)}
    {\(0\)}
    {\(1\)}
\end{problem}

\begin{answer}
A
\end{answer}

\begin{solution}
由 \(\sin x+\cos x=\sqrt2\sin\left(x+\frac\pi4\right)\)，可知 \(-\sqrt2\leq\sin x+\cos x\leq\sqrt2\). 当 \(x+\frac\pi4=\frac{3\pi}{2}+2k\pi\) 时，\(\sin x+\cos x=-\sqrt2\)，下界可以取得. 因而 \(y=\sin x+\cos x+2\) 的最小值为 \(2-\sqrt2\)，故选 A.
\end{solution}

\begin{problem}
设复数 \(z_1=-1-\i\) 在复平面上对应向量 \(\overrightarrow{OZ_1}\)，将 \(\overrightarrow{OZ_1}\) 按顺时针方向旋转 \(\frac{5\pi}{6}\) 后得到向量 \(\overrightarrow{OZ_2}\). 令 \(\overrightarrow{OZ_2}\) 对应的复数为 \(z_2\)，\(z_2\) 的辐角主值为 \(\theta\)，则 \(\tan\theta=\)（\(\quad\)）

\choices
    {\(2 - \sqrt{3}\)}
    {\(-2+\sqrt{3}\)}
    {\(2 + \sqrt{3}\)}
    {\(-2-\sqrt{3}\)}
\end{problem}

\begin{answer}
C
\end{answer}

\begin{solution}
复数 \(z_1=-1-\i\) 的辐角主值为 \(-\frac{3\pi}{4}\). 顺时针旋转 \(\frac{5\pi}{6}\) 后，辐角为 \(-\frac{3\pi}{4}-\frac{5\pi}{6}=-\frac{19\pi}{12}\)，化为辐角主值为 \(\frac{5\pi}{12}\). 所以
\(\tan\theta=\tan\left(\frac\pi4+\frac\pi6\right)=\frac{1+\frac1{\sqrt3}}{1-\frac1{\sqrt3}}=2+\sqrt3\)，故选 C.
\end{solution}

\begin{problem}
设 \(\alpha\)，\(\beta\) 是一个钝角三角形的两个锐角，下列四个不等式中不正确的是（\(\quad\)）

\choices
    {\(\tan \alpha \tan \beta < 1\)}
    {\(\sin \alpha + \sin \beta < \sqrt{2}\)}
    {\(\cos \alpha + \cos \beta > 1\)}
    {\(\frac{1}{2}\tan\left(\alpha+\beta\right)<\tan\frac{\alpha+\beta}{2}\)}
\end{problem}

\begin{answer}
D
\end{answer}

\begin{solution}
因为 \(\alpha\)，\(\beta\) 是钝角三角形的两个锐角，所以 \(0<\alpha+\beta<\frac\pi2\). 令 \(s=\alpha+\beta\). 由正切和角公式，\(\tan s=\frac{\tan\alpha+\tan\beta}{1-\tan\alpha\tan\beta}>0\)，故 \(\tan\alpha\tan\beta<1\)，选项 A 正确. 又
\(\sin\alpha+\sin\beta=2\sin\frac s2\cos\frac{\alpha-\beta}{2}\leq2\sin\frac s2<\sqrt2\)，选项 B 正确.

由于 \(0\leq|\alpha-\beta|<s\)，有
\(\cos\alpha+\cos\beta=2\cos\frac s2\cos\frac{\alpha-\beta}{2}\geq2\cos^2\frac s2=1+\cos s>1\)，选项 C 正确. 最后，令 \(t=\tan\frac s2\)，则 \(0<t<1\)，且 \(\frac12\tan s=\frac{t}{1-t^2}>t=\tan\frac s2\)，所以选项 D 的不等式错误，故选 D.
\end{solution}

\begin{problem}
已知等差数列 \(\{a_n\}\) 满足 \(a_1+a_2+a_3+\cdots+a_{101}=0\)，则有（\(\quad\)）

\choices
    {\(a_{1} + a_{101} > 0\)}
    {\(a_2 + a_{100} < 0\)}
    {\(a_{3} + a_{99} = 0\)}
    {\(a_{51} = 51\)}
\end{problem}

\begin{answer}
C
\end{answer}

\begin{solution}
等差数列的中项为 \(a_{51}\)，且
\(a_1+a_2+\cdots+a_{101}=101a_{51}=0\)，所以 \(a_{51}=0\). 等差数列关于中项对称，因此 \(a_1+a_{101}=2a_{51}=0\)，\(a_2+a_{100}=2a_{51}=0\)，\(a_3+a_{99}=2a_{51}=0\).

所以选项 A 的 \(a_1+a_{101}>0\) 和选项 B 的 \(a_2+a_{100}<0\) 都错误，选项 C 正确；又 \(a_{51}=0\ne51\)，选项 D 错误，故选 C.
\end{solution}

\begin{problem}
已知函数 \(f(x)=ax^3+bx^2+cx+d\) 的图象如图，则（\(\quad\)）

    \bitmapfigure[width=0.34\linewidth]{img/2000/jing_hui_liberal_spring/q14_fig1.png}

\choices
    {\(b\in(-\infty,0)\)}
    {\(b\in(0,1)\)}
    {\(b\in(1,2)\)}
    {\(b\in(2,+\infty)\)}
\end{problem}

\begin{answer}
A
\end{answer}

\begin{solution}
由图象可读出函数的三个零点为 \(0\)，\(1\)，\(2\)，且当 \(x\) 充分大时 \(f(x)>0\)，所以三次项系数 \(a>0\). 因而 \(f(x)=ax(x-1)(x-2)=ax^3-3ax^2+2ax\)，比较 \(x^2\) 项得 \(b=-3a<0\)，故选 A.
\end{solution}

\section{填空题}

\begin{problem}
函数 \(y=\cos\left(\frac{2\pi}{3}x+\frac{\pi}{4}\right)\) 的最小正周期是\fillinblank{}.
\end{problem}

\begin{answer}
\(3\)
\end{answer}

\begin{solution}
函数 \(y=\cos(\omega x+\varphi)\) 的最小正周期为 \(T=\frac{2\pi}{|\omega|}\). 这里 \(\omega=\frac{2\pi}{3}\)，所以 \(T=\frac{2\pi}{2\pi/3}=3\)，填 \(3\).
\end{solution}

\begin{problem}
已知一体积为 \(72\) 的正四面体，连结两个面的重心 \(E\)，\(F\)，则线段 \(EF\) 的长是\fillinblank{}.
\end{problem}

\begin{answer}
\(2\sqrt{2}\)
\end{answer}

\begin{solution}
设正四面体的棱长为 \(\ell\). 取一个面的面积为 \(\frac{\sqrt3}{4}\ell^2\)，正四面体的高为 \(\ell\sqrt{\frac23}\)，于是体积
\(V=\frac13\cdot\frac{\sqrt3}{4}\ell^2\cdot\ell\sqrt{\frac23}=\frac{\ell^3}{6\sqrt2}=72\). 因而 \(\ell^3=432\sqrt2=(6\sqrt2)^3\)，即 \(\ell=6\sqrt2\).

设 \(E\)，\(F\) 分别为两个面的重心. 任取两个面，例如取面 \(ABC\) 和面 \(ABD\)，设四个顶点的位置向量分别为 \(\bs a\)，\(\bs b\)，\(\bs c\)，\(\bs d\)，则
\[
\overrightarrow{EF}=\frac{\bs a+\bs b+\bs d}{3}-\frac{\bs a+\bs b+\bs c}{3}=\frac13\overrightarrow{CD}
\]
故 \(EF=\frac\ell3=2\sqrt2\)，填 \(2\sqrt2\).
\end{solution}

\begin{problem}
\(\left(\sqrt{x}-\frac{1}{\sqrt[3]{x}}\right)^{10}\) 展开式中的常数项是\fillinblank{}.
\end{problem}

\begin{answer}
\(210\)
\end{answer}

\begin{solution}
因为表达式有意义需 \(x>0\)，按二项式定理展开得
\[\left(\sqrt{x}-\frac1{\sqrt[3]{x}}\right)^{10}=\sum_{k=0}^{10}(-1)^k\mathrm C_{10}^k x^{\frac{10-k}{2}-\frac{k}{3}}=\sum_{k=0}^{10}(-1)^k\mathrm C_{10}^k x^{5-\frac56k}
\]
常数项要求指数为 \(0\)，所以 \(5-\frac56k=0\)，得 \(k=6\). 因而常数项为 \((-1)^6\mathrm C_{10}^6=210\)，填 \(210\).
\end{solution}

\begin{problem}
在空间，下列命题正确的是\fillinblank{}.

\begin{circlelist}
\item 如果两直线 \(a\)，\(b\) 分别与直线 \(l\) 平行，那么 \(a \parallel b\)；
\item 如果直线 \(a\) 与平面 \(\beta\) 内的一条直线 \(b\) 平行，那么 \(a \parallel \beta\)；
\item 如果直线 \(a\) 与平面 \(\beta\) 内的两条直线 \(b\)，\(c\) 都垂直，那么 \(a \perp \beta\)；
\item 如果平面 \(\beta\) 内的一条直线 \(a\) 垂直平面 \(\gamma\)，那么 \(\beta \perp \gamma\).
\end{circlelist}
\end{problem}

\begin{answer}
\(\circled{1}\)，\(\circled{4}\)
\end{answer}

\begin{solution}
第一条命题中，直线 \(a\)，\(b\) 都与 \(l\) 平行，因此它们的方向相同，故 \(a\parallel b\)，第一条命题正确.

第二条命题不一定成立. 例如取平面 \(\beta\) 为 \(xy\) 平面，直线 \(b\) 与 \(a\) 都在 \(\beta\) 内且相互平行，则 \(a\parallel b\)，但 \(a\) 在平面 \(\beta\) 内，不是 \(a\parallel\beta\)，所以第二条命题错误.

第三条命题不一定成立，因为平面 \(\beta\) 内的两条直线 \(b\)，\(c\) 可能平行. 此时取平面 \(\beta\) 内一条同时垂直于它们的直线 \(a\)，虽有 \(a\perp b\)，\(a\perp c\)，但 \(a\) 在 \(\beta\) 内，不能推出 \(a\perp\beta\)，所以第三条命题错误.

第四条命题中，直线 \(a\subset\beta\) 且 \(a\perp\gamma\)，说明平面 \(\beta\) 含有平面 \(\gamma\) 的一条法线，因此两平面互相垂直，第四条命题正确. 故应填 \(\circled{1}\)，\(\circled{4}\).
\end{solution}

\section{解答题}

\begin{problem}
已知二次函数 \(f(x)=(\lg a)x^2+2x+4\lg a\) 的最大值为 \(3\)，求 \(a\) 的值.
\end{problem}

\begin{solution}
令 \(t=\lg a\). 因为 \(a>0\)，且二次函数在 \(\R\) 上有最大值，所以 \(t<0\). 配方得

\[
f(x)=t x^2+2x+4t=t\left(x+\frac{1}{t}\right)^2+4t-\frac{1}{t}
\]

因此函数的最大值为 \(4t-\frac{1}{t}\)，由题意

\[
4t-\frac{1}{t}=3
\]

即 \(4t^2-3t-1=0\)，所以 \((4t+1)(t-1)=0\). 由 \(t<0\) 得 \(t=-\frac14\)，从而

\[
a=10^t=10^{-\frac14}
\]
\end{solution}

\begin{problem}
在 \(\triangle ABC\) 中，角 \(A\)，\(B\)，\(C\) 对边分别为 \(a\)，\(b\)，\(c\)，证明：\(\frac{a^2-b^2}{c^2}=\frac{\sin(A-B)}{\sin C}\).
\end{problem}

\begin{solution}
设 \(\triangle ABC\) 的外接圆半径为 \(R\). 由余弦定理，

\(a\cos B=\frac{a^2+c^2-b^2}{2c}\)，\(b\cos A=\frac{b^2+c^2-a^2}{2c}\).

两式相减，得

\[
a\cos B-b\cos A=\frac{a^2-b^2}{c}
\]

所以

\[
a^2-b^2=c\left(a\cos B-b\cos A\right)
\]

由正弦定理 \(a=2R\sin A\)，\(b=2R\sin B\)，于是

\[
a^2-b^2=2Rc\left(\sin A\cos B-\sin B\cos A\right)=2Rc\sin(A-B)
\]

再由正弦定理 \(c=2R\sin C\)，可得

\[
\frac{a^2-b^2}{c^2}=\frac{2R\sin(A-B)}{c}=\frac{\sin(A-B)}{\sin C}
\]
\end{solution}

\begin{problem}
在直角梯形 \(ABCD\) 中，\(\angle D=\angle BAD=90^\circ\)，\(AD=DC=\frac{1}{2}AB=a\)（如图一）. 将 \(\triangle ADC\) 沿 \(AC\) 折起，使 \(D\) 到 \(D'\). 记面 \(ACD'\) 为 \(\alpha\)，面 \(ABC\) 为 \(\beta\)，而面 \(BCD'\) 为 \(\gamma\).

\begin{enumerate}
\item 若二面角 \(\alpha-AC-\beta\) 为直二面角（如图二），求二面角 \(\beta-BC-\gamma\) 的大小；
\item 若二面角 \(\alpha-AC-\beta\) 为 \(60^\circ\)（如图三），求三棱锥 \(D'-ABC\) 的体积.
\end{enumerate}

\FigureLayoutDeclare{img/2000/jing_hui_liberal_spring/q21_fig1.png}{11.0cm}{59bp 22bp 93bp 25bp}
\bitmapfigure[width=0.95\linewidth]{img/2000/jing_hui_liberal_spring/q21_fig1.png}
\end{problem}

\begin{solution}
\begin{enumerate}
\item 取 \(\beta\) 所在平面为 \(z=0\)，建立直角坐标系，使
\(A=(0,0,0)\)，\(B=(2a,0,0)\)，\(D=(0,a,0)\)，\(C=(a,a,0)\).
设二面角 \(\alpha-AC-\beta\) 为 \(\varphi\). 点 \(D\) 在直线 \(AC\) 上的射影为 \(P=(\frac a2,\frac a2,0)\)，且 \(DP=\frac{a}{\sqrt2}\). 因此将 \(PD\) 绕 \(AC\) 旋转 \(\varphi\) 后，
\[
D'=\left(\frac{a(1-\cos\varphi)}2,\frac{a(1+\cos\varphi)}2,\frac{a\sin\varphi}{\sqrt2}\right)
\]

当 \(\varphi=90^{\circ}\) 时，\(D'=\left(\frac a2,\frac a2,\frac a{\sqrt2}\right)\). 令 \(\bs d=\overrightarrow{BC}=(-a,a,0)\)，则平面 \(\beta\) 内垂直于 \(BC\) 的向量可取 \(\bs u=\overrightarrow{CA}=(-a,-a,0)\). 平面 \(\gamma\) 内垂直于 \(BC\) 的向量是 \(\overrightarrow{CD'}\) 在 \(BC\) 的垂线方向上的分量，即
\[
\bs v=\overrightarrow{CD'}-\frac{\overrightarrow{CD'}\cdot\bs d}{\bs d\cdot\bs d}\bs d=\left(-\frac a2,-\frac a2,\frac a{\sqrt2}\right)
\]

设二面角 \(\beta-BC-\gamma\) 为 \(\theta\)，则
\[
\cos\theta=\frac{\bs u\cdot\bs v}{|\bs u|\,|\bs v|}=\frac{a^2}{a\sqrt2\cdot a}=\frac{\sqrt2}{2}
\]
故 \(\theta=45^{\circ}\).

\item 当 \(\varphi=60^{\circ}\) 时，点 \(D'\) 到平面 \(ABC\) 的距离为其纵坐标
\[
h=\frac{a\sin60^{\circ}}{\sqrt2}=\frac{\sqrt6}{4}a
\]
在平面 \(ABC\) 内，\(AB=2a\)，点 \(C\) 到直线 \(AB\) 的距离为 \(a\)，所以
\[
S_{\triangle ABC}=\frac12\cdot2a\cdot a=a^2
\]
因此三棱锥 \(D'-ABC\) 的体积为
\[
V=\frac13S_{\triangle ABC}h=\frac13\cdot a^2\cdot\frac{\sqrt6}{4}a=\frac{\sqrt6}{12}a^3
\]
\end{enumerate}
\end{solution}

\FloatBarrier

\begin{problem}
已知等差数列 \(\{a_n\}\) 的公差和等比数列 \(\{b_n\}\) 的公比相等，且都等于 \(d\)（\(d>0\)，\(d\neq1\)），若 \(a_1=b_1\)，\(a_3=3b_3\)，\(a_5=5b_5\)，求 \(a_n\)，\(b_n\).
\end{problem}

\begin{solution}
设 \(a_1=b_1=A\). 由两数列的公差、公比都为 \(d\)，有
\(a_n=A+(n-1)d\)，\(b_n=A d^{n-1}\).
由已知条件
\(A+2d=3Ad^2\)，\(A+4d=5Ad^4\).
若 \(A=0\)，第一式将推出 \(d=0\)，与 \(d>0\) 矛盾，所以 \(A\neq0\). 由于 \(d>0\)，若 \(3d^2-1=0\) 或 \(5d^4-1=0\)，原来的两式将分别推出 \(0=2d\) 或 \(0=4d\)，均不可能. 因而两式可分别整理为
\(A=\frac{2d}{3d^2-1}\)，\(A=\frac{4d}{5d^4-1}\).
于是
\[
\frac{2d}{3d^2-1}=\frac{4d}{5d^4-1}
\]
因 \(d>0\)，约去 \(2d\)，得
\[
5d^4-1=2\left(3d^2-1\right)
\]
即
\[
5d^4-6d^2+1=0=(5d^2-1)(d^2-1)
\]
由于 \(d>0\) 且 \(d\neq1\)，所以 \(d=\frac1{\sqrt5}\). 代入 \(A+2d=3Ad^2\)，得
\[
A+\frac2{\sqrt5}=\frac35A
\]
从而 \(A=-\sqrt5\). 因此
\(a_n=-\sqrt5+\frac{n-1}{\sqrt5}=\frac{n-6}{\sqrt5}\)，\(b_n=-\sqrt5\left(\frac1{\sqrt5}\right)^{n-1}\).
\end{solution}

\begin{problem}
如图，设点 \(A\) 和 \(B\) 为抛物线 \(y^2=4px\)（\(p>0\)）上原点以外的两个动点，已知 \(OA\perp OB\)，\(OM\perp AB\). 求点 \(M\) 的轨迹方程，并说明它表示什么曲线.

\bitmapfigure[width=0.36\linewidth]{img/2000/jing_hui_liberal_spring/q23_fig1.png}
\end{problem}

\FloatBarrier

\begin{solution}
设抛物线上的点用参数表示为
\(A=(pt^2,2pt)\)，\(B=(pu^2,2pu)\).
由于 \(A\)，\(B\) 都不是原点，故 \(t\neq0\)，\(u\neq0\). 由 \(OA\perp OB\)，得
\[
\overrightarrow{OA}\cdot\overrightarrow{OB}=p^2t^2u^2+4p^2tu=p^2tu\left(tu+4\right)=0
\]
所以 \(tu=-4\). 令 \(s=t+u\)，由两点式可得直线 \(AB\) 的方程
\[
(t+u)y=2x+2ptu
\]
即
\[
2x-sy-8p=0
\]

直线 \(OM\) 垂直于 \(AB\)，而 \((2,-s)\) 是直线 \(AB\) 的一个法向量，所以可设 \(M=\left(2\lambda,-s\lambda\right)\). 将其代入直线 \(AB\) 的方程，得
\[
\lambda\left(s^2+4\right)=8p
\]
从而
\(x=\frac{16p}{s^2+4}\)，\(y=-\frac{8ps}{s^2+4}\).
于是
\[
x^2+y^2=\frac{256p^2+64p^2s^2}{\left(s^2+4\right)^2}=\frac{64p^2}{s^2+4}=4px
\]
所以点 \(M\) 的轨迹满足
\[
x^2+y^2-4px=0
\]
又因为 \(p>0\)，上式参数表示中的 \(x=\frac{16p}{s^2+4}>0\)，故轨迹不包含原点. 反之，圆 \(x^2+y^2=4px\) 上任意非原点的点都有 \(x>0\)，取 \(s=-\frac{2y}{x}\)，则可取 \(t\)，\(u\) 为方程 \(z^2-sz-4=0\) 的两个实根，便能得到对应的点 \(A\)，\(B\). 因此轨迹是
\(x^2+y^2-4px=0\)，\(\left(x\neq0\right)\)
即以 \((2p,0)\) 为圆心、以 \(2p\) 为半径的圆去掉原点.
\end{solution}

\begin{problem}
某地区上年度电价为 \(0.8\text{ 元}/(\mathrm{kW\cdot h})\)，年用电量为 \(a\mathrm{~kW\cdot h}\). 本年度计划将电价降到 \(0.55\text{ 元}/(\mathrm{kW\cdot h})\) 至 \(0.75\text{ 元}/(\mathrm{kW\cdot h})\) 之间，而用户期望电价为 \(0.4\text{ 元}/(\mathrm{kW\cdot h})\). 经测算，下调电价后新增的用电量与实际电价和用户期望电价的差成反比（比例系数为 \(k\)）. 该地区电力的成本为 \(0.3\text{ 元}/(\mathrm{kW\cdot h})\).

\begin{enumerate}
\item 写出本年度电价下调后，电力部门的收益 \(y\) 与实际电价 \(x\) 的函数关系式；
\item 设 \(k = 0.2a\)，当电价最低定为多少时，仍可保证电力部门的收益比上年至少增长 \(20\%\).
\end{enumerate}

注：收益 \(=\) 实际用电量 \(\times\) （实际电价 \(-\) 成本价）.
\end{problem}

\begin{solution}
\begin{enumerate}
\item 设本年度实际电价为 \(x\text{ 元}/(\mathrm{kW\cdot h})\)，则新增用电量为 \(\frac{k}{x-0.4}\mathrm{~kW\cdot h}\)，本年度实际用电量为
\[
a+\frac{k}{x-0.4}
\]
电力部门每用电 \(1\mathrm{~kW\cdot h}\) 的收益为 \(x-0.3\) 元，所以
\(y=\left(a+\frac{k}{x-0.4}\right)(x-0.3)\)，\(0.55\le x\le0.75\).

\item 上年度收益为
\[
a\left(0.8-0.3\right)=0.5a
\]
要比上年至少增长 \(20\%\)，需满足 \(y\ge1.2\times0.5a=0.6a\). 当 \(k=0.2a\) 时，
\[
y=a\left(1+\frac{0.2}{x-0.4}\right)(x-0.3)=a\frac{(x-0.2)(x-0.3)}{x-0.4}
\]
由于 \(0.55\le x\le0.75\)，有 \(x-0.4>0\)，故
\[
\frac{(x-0.2)(x-0.3)}{x-0.4}\ge0.6
\]
等价于
\[
(x-0.2)(x-0.3)\ge0.6(x-0.4)
\]
即
\[
x^2-1.1x+0.3=(x-0.5)(x-0.6)\ge0
\]
在给定范围内 \(x-0.5>0\)，所以必须有 \(x\ge0.6\). 因此电价最低定为 \(0.6\text{ 元}/(\mathrm{kW\cdot h})\).
\end{enumerate}
\end{solution}
